% ============================================================
% Section 102: 氢原子的原子散射因子
% ============================================================
\section{量子力学：氢原子的原子散射因子}
\label{sec:hydrogen-scattering-factor}

\begin{modelbox}[题目：氢原子基态的X射线散射因子]
处于基态的氢原子，电子的概率分布函数为
\begin{equation}
n(r)=(\pi a_0^3)^{-1}e^{-2r/a_0},
\end{equation}
其中 $a_0$ 为玻尔半径。试证明氢原子散射因子为
\begin{equation}
f_G=\frac{16}{(4+G^2a_0^2)^2},
\end{equation}
式中 $G$ 为倒格矢的大小。
\end{modelbox}

\begin{tipbox}[考点快速分析]
\begin{itemize}
    \item \textbf{散射因子定义}：$f_G=\int n(\bm{r})e^{-i\bm{G}\cdot\bm{r}}d^3r$
    \item \textbf{球对称简化}：角向积分给出 $4\pi\int n(r)\dfrac{\sin(Gr)}{Gr}r^2dr$
    \item \textbf{积分技巧}：$\int_0^\infty re^{-\alpha r}\sin(\beta r)dr=\dfrac{2\alpha\beta}{(\alpha^2+\beta^2)^2}$
    \item \textbf{极限检验}：$G=0$ 时 $f_0=1$（氢原子 1 个电子）
\end{itemize}
\end{tipbox}

\begin{derivationbox}[标准解题过程]
\textbf{Step 1: 散射因子的一般表达式}

原子散射因子定义为电子密度的傅里叶变换
\begin{equation}
f_G=\int n(\bm{r})e^{-i\bm{G}\cdot\bm{r}}d^3r.
\end{equation}

对球对称分布，取 $\bm{G}$ 方向为极轴，球坐标中 $\bm{G}\cdot\bm{r}=Gr\cos\theta$，$d^3r=2\pi r^2\sin\theta\,d\theta\,dr$。角向积分
\begin{equation}
\int_0^\pi e^{-iGr\cos\theta}\sin\theta\,d\theta
=\frac{e^{iGr}-e^{-iGr}}{iGr}=\frac{2\sin(Gr)}{Gr}.
\end{equation}

因此
\begin{equation}
f_G=\frac{4\pi}{G}\int_0^\infty rn(r)\sin(Gr)\,dr.
\end{equation}

\textbf{Step 2: 代入氢原子基态密度}

将 $n(r)=(\pi a_0^3)^{-1}e^{-2r/a_0}$ 代入：
\begin{equation}
f_G=\frac{4\pi}{G}\cdot\frac{1}{\pi a_0^3}\int_0^\infty re^{-2r/a_0}\sin(Gr)\,dr
=\frac{4}{a_0^3G}\int_0^\infty re^{-2r/a_0}\sin(Gr)\,dr.
\end{equation}

\textbf{Step 3: 计算径向积分}

令 $\alpha=2/a_0$。利用标准积分公式
\begin{equation}
\int_0^\infty re^{-\alpha r}\sin(\beta r)\,dr=\frac{2\alpha\beta}{(\alpha^2+\beta^2)^2},
\end{equation}
取 $\beta=G$，得
\begin{equation}
\int_0^\infty re^{-2r/a_0}\sin(Gr)\,dr
=\frac{2\cdot(2/a_0)\cdot G}{\bigl[(2/a_0)^2+G^2\bigr]^2}
=\frac{4G/a_0}{\bigl(4/a_0^2+G^2\bigr)^2}
=\frac{4Ga_0^3}{(4+G^2a_0^2)^2}.
\end{equation}

\textbf{Step 4: 得到散射因子}

代回 $f_G$ 表达式：
\begin{equation}
f_G=\frac{4}{a_0^3G}\cdot\frac{4Ga_0^3}{(4+G^2a_0^2)^2}
=\boxed{\frac{16}{(4+G^2a_0^2)^2}.}
\end{equation}

\textbf{检验}：当 $G=0$ 时，$f_0=16/16=1$，恰为氢原子的 1 个电子，符合物理预期。
\end{derivationbox}

\begin{formulabox}[答案汇总]
\begin{equation}
\boxed{f_G=\frac{16}{(4+G^2a_0^2)^2}.}
\end{equation}
\end{formulabox}

\begin{insightbox}[物理图像]
\begin{itemize}
    \item 原子散射因子 $f_G$ 描述原子对 X 射线的散射能力随散射角（即倒格矢 $G$）的变化。$G=0$ 时 $f_0=1$（氢原子），即所有电子同相散射。
    \item $f_G$ 随 $G$ 增大按 $G^{-4}$ 衰减，反映了电子云的空间扩展：电子云越弥散（$a_0$ 越大），$f_G$ 随 $G$ 下降越快。
    \item 氢原子的散射因子是少数可解析求出的原子散射因子之一，常用作更复杂原子散射因子的构成单元或检验理论模型的基准。
\end{itemize}
\end{insightbox}
